\section{Abstract}
\begin{frame}{Abstract}
    \begin{itemize}
        \item \textbf{Problem:} Feed-forward novel view synthesis from \textbf{sparse binocular cameras} is challenging.
        \item \textbf{Limitation of Prior Work:} Gaussian Splatting methods often need dense views or per-scene optimization. Recent feed-forward methods rely on geometry priors requiring large view overlap.
        \item \textbf{Our Solution:} \textbf{Splat-SAP} – a two-stage, feed-forward pipeline for human-centered scenes.
        \begin{itemize}
            \item \textbf{Stage 1:} Reconstruct \textbf{scale-aware point maps} via iterative affinity learning (self-supervised, no 3D data).
            \item \textbf{Stage 2:} Project point maps to target view, refine geometry via stereo matching, and anchor a \textbf{Gaussian plane} for high-quality rendering.
        \end{itemize}
        \item \textbf{Results:} Trained on multi-view human data. Improves point map stability and rendering quality under large sparsity.
        \item Project page: \url{https://yaourtb.github.io/Splat-SAP}
    \end{itemize}
\end{frame}

\section{Introduction}
\begin{frame}{Introduction: Motivation \& Challenges}
    \begin{itemize}
        \item \textbf{Goal:} Real-time, feed-forward free-viewpoint video from sparse views (e.g., telecommunication, broadcast).
        \item \textbf{Gaussian Splatting} is fast but typically requires \textbf{per-scene optimization} and \textbf{dense inputs}.
        \item Recent feed-forward Gaussian methods (e.g., MVSNeRF, GPS-Gaussian) need \textbf{large view overlap} for reliable geometry, failing under \textbf{large sparsity}.
        \item \textbf{DUSt3R} offers a promising \textbf{point map} representation robust to sparsity but is \textbf{scale-invariant}, causing instability in video.
        \item \textbf{Key Gaps:}
        \begin{itemize}
            \item Achieving \textbf{metric-scale} geometry without 3D supervision.
            \item Enforcing \textbf{stereo consistency} for dynamic scenes.
            \item Rendering efficiently from sparse binocular inputs.
        \end{itemize}
    \end{itemize}
\end{frame}

\begin{frame}{Introduction: Our Approach \& Contributions}
    \begin{itemize}
        \item \textbf{Splat-SAP:} A coarse-to-fine pipeline for metric reconstruction and rendering.
        \begin{center}
            \includegraphics[width=0.9\textwidth]{figs/0_teaser.pdf}
        \end{center}
        \item \textbf{Core Ideas:}
        \begin{enumerate}
            \item \textbf{Affinity Learning:} Transform DUSt3R's canonical point maps to real metric space using a learned \textbf{scaling factor} and \textbf{pixel-wise translation}.
            \item \textbf{Gaussian Plane Rendering:} Anchor Gaussians on a target-view plane, initialized via projected point maps and refined via a \textbf{3D cost volume}.
        \end{enumerate}
        \item \textbf{Contributions:}
        \begin{enumerate}
            \item Self-supervised, feed-forward pipeline for scale-aware point maps and rendering.
            \item Coarse (2D affinity) to fine (3D stereo) registration strategy.
            \item Efficient Gaussian plane rendering incorporating 2D/3D features.
        \end{enumerate}
    \end{itemize}
\end{frame}

\section{Related Work}
\begin{frame}{Related Work: Novel View Synthesis}
    \begin{itemize}
        \item \textbf{NeRF-based Methods:} High quality but slow. \textbf{Generalizable NeRFs} (e.g., IBRNet, PixelNeRF) use learned priors for feed-forward inference.
        \item \textbf{3D Gaussian Splatting (3D-GS):} Real-time rendering via differentiable splatting.
        \begin{itemize}
            \item \textbf{Optimization-based:} Require per-scene tuning (minutes).
            \item \textbf{Feed-forward:} Predict Gaussians from images using geometry proxies (MVS, stereo). \textbf{Still need good view overlap}.
        \end{itemize}
        \item \textbf{Our Position:} Use a robust point map prior (DUSt3R) and refine it for sparse, wide-baseline rendering.
    \end{itemize}
\end{frame}

\begin{frame}{Related Work: 3D Reconstruction}
    \begin{itemize}
        \item \textbf{Traditional MVS \& Stereo:} Require dense views and overlap. Struggle with sparsity.
        \item \textbf{Neural Implicit Surfaces:} (e.g., NeuS) reconstruct from images via rendering loss, but need optimization.
        \item \textbf{Point Map Representations (DUSt3R):} Predict a 3D point per pixel from 2+ views, \textbf{without camera poses}. Robust to sparsity but \textbf{scale-invariant}.
        \item \textbf{Follow-ups:} MASt3R, VGGT improve but remain scale-free. Methods like NoPoSplat use it for rendering but lack stereo constraint for dynamics.
        \item \textbf{Our Position:} Add metric scale and stereo consistency to the point map foundation for stable dynamic scene rendering.
    \end{itemize}
\end{frame}

\section{Method}
\begin{frame}{Method: Pipeline Overview}
    \begin{center}
        \includegraphics[width=0.95\textwidth]{figs/pipeline.pdf} % Assuming a pipeline figure exists
    \end{center}
    \begin{itemize}
        \item \textbf{Input:} Two sparse source views ($I^l$, $I^r$) + camera calibration.
        \item \textbf{Stage 1 - Geometry:} Predict scale-aware point maps $X_t^i$ via affinity learning.
        \item \textbf{Stage 2 - Rendering:}
        \begin{enumerate}
            \item Project $X_t^i$ to target view for initial depth.
            \item Refine depth via 3D cost volume (stereo matching).
            \item Build \& splat a \textbf{Gaussian plane} $\mathcal{G}$ on target view.
        \end{enumerate}
        \item \textbf{Output:} High-quality novel view image.
    \end{itemize}
\end{frame}

\begin{frame}{Method: Stage 1 - Scale-Aware Point Maps}
    \begin{itemize}
        \item \textbf{Base Representation:} Use \textbf{MASt3R} to get canonical point maps $X^i$ (scale-invariant).
        \item \textbf{Affine Transformation to Real Space:} $X_t^i = S X^i + T^i$
        \begin{itemize}
            \item \textbf{Scaling Factor $S$:} Global, from camera intrinsics ($f$) and baseline ($d$).
            \begin{itemize}
                \item Encode $f, d$. Use self/cross-attention on ViT features to get global context.
                \item $S = MLP(\text{global features}, \text{encoding})$.
            \end{itemize}
            \item \textbf{Translation Map $T^i$:} Pixel-wise, corrects misalignment between views.
            \begin{itemize}
                \item Project rescaled points $SX^i$ to other view, query features.
                \item Iteratively update $T^i$ with a GRU using features from both views.
            \end{itemize}
        \end{itemize}
        \item \textbf{Self-Supervised Training:} No 3D ground truth needed.
        \begin{itemize}
            \item Build auxiliary Gaussian planes on source views, render target view.
            \item Loss: $\mathcal{L}_{stage1} = \mathcal{L}_{render} + \gamma \mathcal{L}_{CD}$ (Chamfer Distance between point sets).
        \end{itemize}
    \end{itemize}
\end{frame}

\begin{frame}{Method: Stage 2 - Rendering via Gaussian Plane}
    \begin{itemize}
        \item \textbf{Step 1: 3D Geometry Refinement}
        \begin{itemize}
            \item Encode source images in fine resolution.
            \item Get initial target depth $\mathcal{D}^k$ by $\alpha$-blending projected $X_t^i$.
            \item For each pixel: sample depth candidates, warp features from both views, build 3D cost volume $\Phi^k$.
            \item Regress refined depth $\bar{d}$ from probability distribution.
        \end{itemize}
        \item \textbf{Step 2: Gaussian Plane $\mathcal{G}$ Construction}
        \begin{itemize}
            \item \textbf{Position:} From refined depth $\bar{d}$.
            \item \textbf{Color:} Initialize by blending warped source colors, then add learned residual.
            \item \textbf{Attributes ($\mathcal{P}_r, \mathcal{P}_s, \mathcal{P}_o$):} Predict from aggregated 2D/3D feature map $\mathcal{M}$.
            \item $\mathcal{G} = \{\mathcal{P}_c, \mathcal{P}_r, \mathcal{P}_s, \mathcal{P}_o\}$ is a 2D map of Gaussians anchored on the target view.
        \end{itemize}
        \item \textbf{Step 3: Splatting \& Loss}
        \begin{itemize}
            \item Splat $\mathcal{G}$ to render high-res image.
            \item Loss: $\mathcal{L}_{stage2} = \lambda_1 \mathcal{L}_{render}(\hat{I}_f) + \lambda_2 \mathcal{L}_{render}(\hat{I}_h)$.
        \end{itemize}
    \end{itemize}
\end{frame}

\section{Experiments}
\begin{frame}{Experiments: Settings}
    \begin{itemize}
        \item \textbf{Datasets:} Multi-view human scenes from 3 camera types.
        \begin{itemize}
            \item \textbf{Industry Camera (THumanMV):} 15 train / 11 val sequences.
            \item \textbf{Mobile Phone (4K4D, SelfCap):} $\sim$3000 train frames.
            \item \textbf{GoPro (Our Capture):} 6 train / 4 val sequences (sports, multi-person).
        \end{itemize}
        \item \textbf{Setup:} Feed leftmost \& rightmost cameras (large baseline). Render 4 middle views.
        \item \textbf{Metrics:}
        \begin{itemize}
            \item \textbf{Rendering:} PSNR, SSIM, LPIPS.
            \item \textbf{Geometry:} Chamfer Distance (CD) to SfM point cloud.
        \end{itemize}
        \item \textbf{Baselines:}
        \begin{itemize}
            \item \textbf{Rendering:} ENeRF, MVSplat, MVSGaussian, NoPoSplat, 4D-GS.
            \item \textbf{Geometry:} DUSt3R, MASt3R, VGGT, Pow3R, Prompt-DA.
        \end{itemize}
    \end{itemize}
\end{frame}

\begin{frame}{Experiments: Rendering Results}
    \begin{itemize}
        \item \textbf{Quantitative:} Our method outperforms baselines, especially on Camera and GoPro data.
        \begin{center}
            \includegraphics[width=0.8\textwidth]{tables/quantitative.png} % Placeholder for table
        \end{center}
        \item \textbf{Qualitative:}
        \begin{center}
            \includegraphics[width=0.9\textwidth]{figs/4_comp.pdf}
        \end{center}
        \begin{itemize}
            \item MVSGaussian/ENeRF: Missing thin structures due to sparsity.
            \item MVSplat/NoPoSplat: Misaligned Gaussians, artifacts.
            \item Ours: Higher quality, better details and consistency.
        \end{itemize}
    \end{itemize}
\end{frame}

\begin{frame}{Experiments: Geometry \& Ablation Results}
    \begin{itemize}
        \item \textbf{Geometry Comparison:} Our affinity learning achieves best Chamfer Distance.
        \begin{center}
            \includegraphics[width=0.7\textwidth]{tables/geometry.png} % Placeholder
        \end{center}
        \item \textbf{Ablation - Affinity:}
        \begin{center}
            \includegraphics[width=0.9\textwidth]{figs/4_abla_geo.pdf}
        \end{center}
        \begin{itemize}
            \item Scaling only ($S$): Misalignment remains.
            \item Full ($S+T$): Corrects alignment, better CD.
        \end{itemize}
        \item \textbf{Ablation - Rendering:}
        \begin{center}
            \includegraphics[width=0.9\textwidth]{figs/4_abla_full.pdf}
        \end{center}
        \begin{itemize}
            \item Stage 1 only: Holes, artifacts.
            \item Stage 2 with refinement: Fills holes, refines details.
        \end{itemize}
    \end{itemize}
\end{frame}

\section{Conclusion}
\begin{frame}{Conclusion \& Future Work}
    \begin{itemize}
        \item \textbf{Summary:}
        \begin{itemize}
            \item Proposed \textbf{Splat-SAP}, a feed-forward pipeline for sparse-view human scene synthesis.
            \item Introduced self-supervised \textbf{scale-aware point map} reconstruction via affinity learning.
            \item Designed a \textbf{Gaussian plane} rendered with coarse-to-fine geometry.
            \item Demonstrated superior rendering and geometry quality vs. SOTA.
        \end{itemize}
        \item \textbf{Limitations \& Future Work:}
        \begin{itemize}
            \item \textbf{Floating artifacts} from MASt3R's uncertain boundary points.
            \item Future: Incorporate monocular depth priors to clean up artifacts.
            \item Extend to more extreme sparsity or moving cameras.
        \end{itemize}
    \end{itemize}
\end{frame}

\section{Appendix}
\begin{frame}{Appendix: More Geometry Results}
    \begin{center}
        \includegraphics[width=0.95\textwidth]{figs/x_comp_geo.pdf}
    \end{center}
    \begin{itemize}
        \item Our method maintains better inter-view consistency for the foreground human compared to DUSt3R, Pow3R, VGGT, and Prompt-DA.
    \end{itemize}
\end{frame}

\begin{frame}{Appendix: Training \& Time Analysis}
    \begin{itemize}
        \item \textbf{Stage 1 Training:} Self-supervised via rendering and Chamfer loss.
        \begin{center}
            \includegraphics[width=0.6\textwidth]{figs/train_stage1_aaai65.pdf}
        \end{center}
        \item \textbf{Time Analysis (RTX 3090):}
        \begin{center}
            \includegraphics[width=0.7\textwidth]{tables/time_analysis.png} % Placeholder
        \end{center}
        \begin{itemize}
            \item MASt3R inference is the bottleneck.
            \item Our affinity and rendering modules are efficient.
            \item Potential for further optimization (e.g., TensorRT).
        \end{itemize}
    \end{itemize}
\end{frame}